\documentclass[11pt]{article}

\usepackage{amsmath,amssymb,amsthm}

\newtheorem{theorem}{Theorem}
\newtheorem{lemma}{Lemma}

\newcommand{\D}{\mathbb D}
\newcommand{\C}{\mathbb C}
\newcommand{\alg}{\operatorname{alg}}
\newcommand{\RePart}{\operatorname{Re}}

\begin{document}

\begin{theorem}
Let \(A\in\C^{n\times n}\), and let
\[
 W(A)=\{x^*Ax:x\in\C^n,\ \|x\|_2=1\}.
\]
Then every complex polynomial \(p\) satisfies
\[
 \|p(A)\|_2\leq 2\max_{z\in W(A)}|p(z)|.
\]
\end{theorem}

We write \(\RePart X=(X+X^*)/2\), and \(\alg(T)\) for the unital
algebra generated by a matrix \(T\).  The main point of the proof is the
following positive-real completion lemma.

\begin{lemma}[Positive-real completion]
\label{lem:completion}
Suppose that \(T\in\C^{n\times n}\) has \(n\) distinct eigenvalues, all in
the closed unit disk.  Suppose that \(H:\D\to\C^{n\times n}\) is analytic
and has the following properties:
\[
 \begin{gathered}
 H(0)=I,\qquad \RePart H(z)\succeq0\quad(z\in\D),\\
 H(z)-(I-zT)^{-1}\in\alg(T^*)\quad(z\in\D).
 \end{gathered}
\]
Then \(\|T\|_2\leq2\).
\end{lemma}

\begin{proof}
We first recall a consequence of the matrix Herglotz theorem.  If
\(K:\D\to\C^{n\times n}\) is analytic and \(\RePart K\succeq0\), then
\begin{equation}
 \mathcal L_K(z,w)=\frac{K(z)+K(w)^*}{1-z\overline w}
 \label{eq:herglotz-kernel}
\end{equation}
is a positive matrix-valued kernel.  In other words, for arbitrary points
\(z_0,\ldots,z_N\in\D\) and vectors
\(\xi_0,\ldots,\xi_N\in\C^n\),
\begin{equation}
 \sum_{i,j=0}^N
 \xi_i^*\mathcal L_K(z_i,z_j)\xi_j\geq0.
 \label{eq:kernel-positive}
\end{equation}
For completeness, the matrix Herglotz representation is obtained by
applying the scalar Herglotz representation to \(x^*K(z)x\) and using
polarization.  It gives a positive matrix-valued measure \(M\) on the unit
circle and a self-adjoint matrix \(C\) such that
\[
 K(z)=iC+\int_{|\zeta|=1}\frac{\zeta+z}{\zeta-z}\,dM(\zeta).
\]
The elementary identity
\[
 \frac{1}{1-z\overline w}
 \left(
 \frac{\zeta+z}{\zeta-z}
 +
 \overline{\frac{\zeta+w}{\zeta-w}}
 \right)
 =
 \frac{2}{(1-z\overline\zeta)(1-\overline w\zeta)}
\]
proves \eqref{eq:kernel-positive}.

Choose a diagonalization
\[
 T=S\Lambda S^{-1},
 \qquad
 \Lambda=\operatorname{diag}(\lambda_1,\ldots,\lambda_n),
 \qquad
 G=S^*S.
\]
Every element of \(\alg(T^*)\) has the form \(S^{-*}DS^*\), with \(D\)
diagonal.  The map \(D\mapsto S^{-*}DS^*\) is a fixed linear isomorphism
from the diagonal matrices onto \(\alg(T^*)\).  Thus, when its inverse is
applied to the analytic function
\(H(z)-(I-zT)^{-1}\), the resulting diagonal matrix \(D(z)\) is analytic
and satisfies \(D(0)=0\).  Consequently, after putting
\(K(z)=S^*H(z)S\), the hypotheses give
\begin{equation}
 K(z)=G(I-z\Lambda)^{-1}+D(z)G,
 \label{eq:K-form}
\end{equation}
where \(D(z)=\operatorname{diag}(d_1(z),\ldots,d_n(z))\).  Congruence by
\(S\) preserves positive real parts, so the kernel in
\eqref{eq:herglotz-kernel}, formed from this \(K\), is positive.

Define three Hermitian matrices by
\begin{equation}
 \begin{split}
 P_{ij}&=\frac{G_{ij}}{1-\overline{\lambda_i}\lambda_j/4},\\
 R_{ij}&=\frac{G_{ij}}{1-\overline{\lambda_i}\lambda_j/2},\\
 X&=R-P.
 \end{split}
 \label{eq:PRX}
\end{equation}
Fix \(u=(u_1,\ldots,u_n)^T\in\C^n\).  In
\eqref{eq:kernel-positive}, use the \(n\) points and vectors
\[
 z_i=\frac{\overline{\lambda_i}}2,
 \qquad
 \xi_i=u_i e_i
 \quad(1\leq i\leq n),
\]
and also use the point \(z_0=0\) with the vector
\begin{equation}
 \xi_0=v=-G^{-1}Pu.
 \label{eq:v-choice}
\end{equation}
This choice removes the entire unknown term \(D\) in \eqref{eq:K-form}.
Indeed, its contribution to the quadratic form in
\eqref{eq:kernel-positive} is
\[
 2\RePart\sum_{i=1}^n
 \overline{u_i}d_i(z_i)
 \left(
 (Gv)_i+
 \sum_{j=1}^n
 \frac{G_{ij}u_j}{1-z_i\overline{z_j}}
 \right).
\]
The sum in parentheses is \((Gv+Pu)_i\), which is zero by
\eqref{eq:v-choice}.

It remains to calculate the contribution of
\(G(I-z\Lambda)^{-1}\).  Between the points \(z_i,z_j\) it is
\[
 \frac{2G_{ij}}
 {(1-\overline{\lambda_i}\lambda_j/2)
  (1-\overline{\lambda_i}\lambda_j/4)}
 =(4R-2P)_{ij}.
\]
Between \(z_i\) and \(0\) it is \((G+R)_{ij}\), and at \(0,0\) it is
\(2G\).  Hence \eqref{eq:kernel-positive} and \eqref{eq:v-choice} give,
for every \(u\),
\[
 \begin{split}
 0\leq{}&u^*(4R-2P)u
       +2\RePart\bigl(u^*(G+R)v\bigr)+2v^*Gv\\
 ={}&u^*\bigl(4X-XG^{-1}P-PG^{-1}X\bigr)u.
 \end{split}
\]
Thus
\begin{equation}
 4X-XG^{-1}P-PG^{-1}X\succeq0.
 \label{eq:X-inequality}
\end{equation}

Set
\[
 \begin{aligned}
 C&=G^{1/2}\Lambda G^{-1/2},
 &\widehat P&=G^{-1/2}PG^{-1/2},\\
 \widehat R&=G^{-1/2}RG^{-1/2},
 &\widehat X&=\widehat R-\widehat P.
 \end{aligned}
\]
Expanding the two scalar kernels in \eqref{eq:PRX} as geometric series
gives
\begin{equation}
 \widehat P=\sum_{k=0}^{\infty}4^{-k}C^{*k}C^k,
 \qquad
 \widehat R=\sum_{k=0}^{\infty}2^{-k}C^{*k}C^k,
 \label{eq:gramians}
\end{equation}
and therefore
\begin{equation}
 \widehat X=
 \sum_{k=1}^{\infty}(2^{-k}-4^{-k})C^{*k}C^k\succeq0.
 \label{eq:X-positive}
\end{equation}
Indeed,
\[
 C^k=G^{1/2}\Lambda^kG^{-1/2},
\]
so the powers \(C^k\) are uniformly bounded because
\(|\lambda_i|\leq1\).  The geometric weights therefore make both series
converge in norm.  Congruencing \eqref{eq:X-inequality} by \(G^{-1/2}\)
yields
\begin{equation}
 4\widehat X-\widehat X\widehat P-
 \widehat P\widehat X\succeq0.
 \label{eq:hat-inequality}
\end{equation}

Let \(e\) be an eigenvector of the Hermitian matrix \(\widehat P\), with
\(\widehat Pe=pe\).  If \(p>2\), then \eqref{eq:hat-inequality} gives
\[
 0\leq 2(2-p)e^*\widehat Xe.
\]
Since \(\widehat X\succeq0\), this forces \(\widehat Xe=0\).  The
\(k=1\) summand in \eqref{eq:X-positive} is \(C^*C/4\), so \(Ce=0\).
Formula \eqref{eq:gramians} then gives \(\widehat Pe=e\), contradicting
\(p>2\).  It follows that
\begin{equation}
 I\preceq\widehat P\preceq2I.
 \label{eq:P-bounds}
\end{equation}
The first series in \eqref{eq:gramians} satisfies the Stein identity
\[
 \widehat P=I+\frac14 C^*\widehat PC.
\]
Using \eqref{eq:P-bounds}, we conclude that
\[
 C^*C\preceq C^*\widehat PC
 =4(\widehat P-I)\preceq4I.
\]
Thus \(\|C\|_2\leq2\).  Finally, if \(S=UG^{1/2}\) is the polar
decomposition of \(S\), then
\[
 T=UCU^*.
\]
Therefore \(\|T\|_2\leq2\), as claimed.
\end{proof}

We next construct the positive-real completion from the numerical range.

\begin{lemma}[The double-layer positive map]
\label{lem:double-layer}
Let \(B\in\C^{n\times n}\), and let \(\Omega\) be a bounded convex domain
with continuously differentiable boundary such that
\(W(B)\subset\Omega\).  If \(f\) is analytic on a neighborhood of
\(\overline\Omega\) and \(\max_{\overline\Omega}|f|\leq1\), put
\(T=f(B)\).  If \(\alg(T)=\alg(B)\), then there is an analytic function
\(H:\D\to\C^{n\times n}\) satisfying all the hypotheses of
Lemma~\ref{lem:completion}, except possibly the distinct-eigenvalue
hypothesis on \(T\).
\end{lemma}

\begin{proof}
We also have
\[
 \sigma(B)\subset W(B)\subset\Omega.
\]
Thus the spectral mapping theorem and the bound on \(f\) give
\[
 \sigma(T)=f(\sigma(B))\subset\overline{\D}.
\]

Orient \(\partial\Omega\) counterclockwise.  Let \(\nu(\sigma)\) denote
the outward unit normal and let \(ds\) denote arclength.  Put
\[
 R_\sigma=(\sigma I-B)^{-1}
\]
and define an operator-valued measure on \(\partial\Omega\) by
\begin{equation}
 d\mu(\sigma)=\frac1{2\pi}
 \left(\nu(\sigma)R_\sigma+
 \overline{\nu(\sigma)}R_\sigma^*\right)ds.
 \label{eq:mu-definition}
\end{equation}
This measure is positive.  Indeed,
\[
 \nu(\sigma)(\overline\sigma I-B^*)+
 \overline{\nu(\sigma)}(\sigma I-B)
 =2\RePart\bigl(\overline{\nu(\sigma)}(\sigma I-B)\bigr)
 \succeq0,
\]
because the tangent line at \(\sigma\) supports the convex set
\(\overline\Omega\) and \(W(B)\subset\Omega\).  Multiplication on the left
by \(R_\sigma^*\) and on the right by \(R_\sigma\) gives precisely the
matrix in parentheses in \eqref{eq:mu-definition}.

Along the positively oriented boundary, \(d\sigma=i\nu(\sigma)ds\).
Cauchy's formula therefore gives
\[
 \frac1{2\pi}\int_{\partial\Omega}
 \nu(\sigma)R_\sigma\,ds=I.
\]
Taking adjoints deals with the second term in
\eqref{eq:mu-definition}, and hence
\begin{equation}
 \int_{\partial\Omega}d\mu=2I.
 \label{eq:mu-mass}
\end{equation}
Consequently,
\begin{equation}
 \Phi(h)=\frac12\int_{\partial\Omega}h(\sigma)\,d\mu(\sigma)
 \label{eq:Phi}
\end{equation}
is a unital positive map from \(C(\partial\Omega)\) into
\(\C^{n\times n}\).

For a function \(h\) analytic near \(\overline\Omega\), define
\begin{equation}
 g_h(z)=\frac1{2\pi i}\int_{\partial\Omega}
 \frac{\overline{h(\sigma)}}{\sigma-z}\,d\sigma,
 \qquad z\in\Omega.
 \label{eq:companion}
\end{equation}
The first part of \eqref{eq:mu-definition}, Cauchy's formula, and the
adjoint of \eqref{eq:companion} give
\begin{equation}
 2\Phi(h)=h(B)+g_h(B)^*.
 \label{eq:Phi-identity}
\end{equation}

For \(z\in\D\), consider the boundary function
\[
 h_z(\sigma)=\frac{1+zf(\sigma)}{1-zf(\sigma)}.
\]
It has nonnegative real part.  Define
\[
 H(z)=\Phi(h_z).
\]
The map \(H\) is analytic and \(H(0)=I\).  Since the measure \(\mu\) is
Hermitian, \(\Phi(\overline h)=\Phi(h)^*\); hence
\[
 \RePart H(z)=\Phi(\RePart h_z)\succeq0.
\]
Since
\[
 h_z=1+2\sum_{m=1}^{\infty}z^m f^m,
\]
where the series converges uniformly on \(\partial\Omega\) for each
fixed \(z\in\D\), the boundedness of \(\Phi\) and
\eqref{eq:Phi-identity} show that, with norm convergence,
\begin{equation}
 H(z)=I+\sum_{m=1}^{\infty}z^m
 \left(T^m+g_{f^m}(B)^*\right).
 \label{eq:H-expansion}
\end{equation}
It follows that
\begin{equation}
 H(z)-(I-zT)^{-1}
 =\sum_{m=1}^{\infty}z^m g_{f^m}(B)^*.
 \label{eq:completion}
\end{equation}
Every \(g_{f^m}(B)\) belongs to \(\alg(B)\), by the finite-dimensional
holomorphic functional calculus (equivalently, by Hermite
interpolation).  The assumption
\(\alg(T)=\alg(B)\) therefore puts every partial sum on the right-hand
side of \eqref{eq:completion} in \(\alg(T^*)\).  The series converges in
norm: the series in \eqref{eq:H-expansion} does, as just noted, and
\(\sum_{m\geq0}z^mT^m=(I-zT)^{-1}\) converges because
\(\rho(zT)<1\).  Finally, \(\alg(T^*)\) is finite-dimensional and hence
norm-closed, so the limit in \eqref{eq:completion} also belongs to
\(\alg(T^*)\).  This proves the lemma.
\end{proof}

\begin{proof}[Proof of the theorem]
The numerical range is compact and convex.  First choose a bounded convex
domain \(\Omega\), with smooth boundary, such that
\[
 W(A)\subset\Omega.
\]
We shall prove
\begin{equation}
 \|p(A)\|_2\leq2\max_{z\in\overline\Omega}|p(z)|.
 \label{eq:Omega-bound}
\end{equation}

Matrices with distinct eigenvalues are dense in
\(\C^{n\times n}\).  Moreover,
\[
 W(B)\subset W(A)+\{z:|z|\leq\|B-A\|_2\}.
\]
Thus \(A\) can be approximated by matrices \(B\) with distinct
eigenvalues and with \(W(B)\subset\Omega\).  Since polynomial functional
calculus is continuous, it is enough to establish
\eqref{eq:Omega-bound} for each such \(B\).

Fix one such \(B\), and write its distinct eigenvalues as
\(\beta_1,\ldots,\beta_n\).  Put
\[
 M_\Omega=\max_{z\in\overline\Omega}|p(z)|.
\]
If \(M_\Omega=0\), then \(p\) is identically zero and there is nothing to
prove.  Otherwise, put \(f=p/M_\Omega\), and set
\[
 R_\Omega=\max_{z\in\overline\Omega}|z|.
\]
For a complex number \(\eta\) close to zero, define
\begin{equation}
 f_\eta(z)=\frac{f(z)+\eta z}{1+|\eta|R_\Omega}.
 \label{eq:f-eta}
\end{equation}
Then \(\max_{\overline\Omega}|f_\eta|\leq1\).  Apart from finitely many
values of \(\eta\), the numbers
\[
 f(\beta_i)+\eta\beta_i,
 \qquad 1\leq i\leq n,
\]
are pairwise distinct.  Choose a sequence of nonzero admissible values
of \(\eta\) tending to zero.  For each of them, the matrix
\[
 T_\eta=f_\eta(B)
\]
has distinct eigenvalues.  It also generates the same algebra as \(B\).
Indeed, \(T_\eta\in\alg(B)\), while Lagrange interpolation at the distinct
numbers \(f_\eta(\beta_i)\) gives a polynomial \(q\) such that
\(q(T_\eta)=B\).  Hence
\[
 \alg(T_\eta)=\alg(B).
\]

Apply Lemma~\ref{lem:double-layer} to \(B\) and \(f_\eta\).  The resulting
positive-real function satisfies the hypotheses of
Lemma~\ref{lem:completion}; indeed, spectral mapping and
\(\max_{\overline\Omega}|f_\eta|\leq1\) put every eigenvalue of
\(T_\eta\) in the closed unit disk.  Hence
\[
 \|f_\eta(B)\|_2\leq2.
\]
Letting \(\eta\to0\) gives \(\|f(B)\|_2\leq2\), or
\[
 \|p(B)\|_2\leq2M_\Omega.
\]
Now let the simple-spectrum matrices \(B\) tend to \(A\).  This proves
\eqref{eq:Omega-bound}.

Finally, a compact convex set in the plane admits smooth bounded convex
outer approximations converging to it in Hausdorff distance.  Choose such
domains \(\Omega_k\) for \(W(A)\).  Uniform continuity of \(p\) on a fixed
compact neighborhood of \(W(A)\) gives
\[
 \max_{z\in\overline\Omega_k}|p(z)|
 \longrightarrow
 \max_{z\in W(A)}|p(z)|.
\]
Applying \eqref{eq:Omega-bound} to \(\Omega_k\) and taking the limit proves
\[
 \|p(A)\|_2\leq2\max_{z\in W(A)}|p(z)|.
\]
Since \(W(A)\) is already compact, this is equivalently the assertion
that its closure is a \(2\)-spectral set for \(A\).  To spell out the
standard equivalence, the spectrum of \(A\) lies in \(W(A)\).  If a
rational function \(r\) has no pole on the compact convex set \(W(A)\),
choose \(\delta>0\) so that it has no pole on the closed
\(\delta\)-neighborhood of \(W(A)\).  That neighborhood is convex, so
Runge approximation gives polynomials \(q_j\) converging uniformly to
\(r\) there.  Hence \(q_j(A)\to r(A)\) by the holomorphic functional
calculus, while
\(\max_{W(A)}|q_j|\to\max_{W(A)}|r|\).  Applying the polynomial estimate
to \(q_j\) and passing to the limit gives the same estimate for \(r\),
which is precisely the \(2\)-spectral-set assertion.
\end{proof}

\end{document}
